
%%%%%%%%%%%%%%%%%%%%%%%%%%%%%%%%%%%%%%%%%
%
% Copyright Piotr Sroczkowski
%
%%%%%%%%%%%%%%%%%%%%%%%%%%%%%%%%%%%%%%%%%

%----------------------------------------------------------------------------------------
%	PACKAGES AND DOCUMENT CONFIGURATIONS
%----------------------------------------------------------------------------------------

\documentclass{article}

\usepackage[version=3]{mhchem} % Package for chemical equation typesetting
\usepackage{siunitx} % Provides the \SI{}{} and \si{} command for typesetting SI units
\usepackage{graphicx} % Required for the inclusion of images
\usepackage{natbib} % Required to change bibliography style to APA
\usepackage{amsmath} % Required for some math elements 
\usepackage{datetime}
\usepackage[utf8]{inputenc}
\usepackage[polish]{babel}
\usepackage{polski}


\setlength\parindent{0pt} % Removes all indentation from paragraphs

\renewcommand{\labelenumi}{\alph{enumi}.} % Make numbering in the enumerate environment by letter rather than number (e.g. section 6)

%\usepackage{times} % Uncomment to use the Times New Roman font

%----------------------------------------------------------------------------------------
%	DOCUMENT INFORMATION
%----------------------------------------------------------------------------------------

\title{Przedmiot specjalnościowy\\ Programowanie interfejsu użytkownika internetu\\ Edytor } % Title


\author{Piotr \textsc{Sroczkowski}} % Author name

\date{\today} % Date for the report

\begin{document}

\maketitle % Insert the title, author and date


\begin{center}
\includegraphics{logo.png}
\begin{tabular}{l r}
Prowadzący: & dr inż Marek Żabka % Instructor/supervisor
\end{tabular}
\end{center}

\pagebreak

% If you wish to include an abstract, uncomment the lines below
% \begin{abstract}
% Abstract text
% \end{abstract}

%----------------------------------------------------------------------------------------
%	SECTION 1
%----------------------------------------------------------------------------------------


\section{Założenia projektu strony}
\begin{enumerate}
    \item Dodawanie elementu do diagramu
    \item Przesuwanie elementu
    \item Sprawdzanie kolizji z innymi elementami i granicami daigramu przy przesuwaniu elementu
    \item Sprawdzanie kolizji przy dodawaniu nowego elementu
    \item Zmiana tekstu elementu
    \item Łączenie elementów
\end{enumerate}

\end{document}
